
\documentclass[11pt]{article}
\usepackage[margin=1in]{geometry}

% Core packages
\usepackage{amsmath,amssymb}
\usepackage[spanish]{babel}
\usepackage{tikz-cd}
\usepackage{physics}
\usepackage{bm}
\usepackage{dsfont}
\usepackage{multicol}
\usepackage{tikz}
\usetikzlibrary{decorations.pathmorphing, patterns}

% Paragraphs
\setlength{\parindent}{0pt}
\setlength{\parskip}{1\baselineskip}

\title{Ayudantía MMF II}
\author{Jovanny Jara}
\date{10 de abril de 2026}

\begin{document}
\maketitle

\section{Determinantes de orden $n$}

Para encontrar el determinante de una matriz cuadrada $n\times n$ o de orden $n$, podemos usar la expansión por menores, o la regla de Laplace, para hacernos mas fácil el trabajo.

Los menores, denotados por $M_{ij}$, son los determinantes correspondientes al elemento ($ij$) de la matriz original, denotado por $a_{ij}$, siendo estos obtenidos de eliminar la fila $i$ y la columna $j$ de la matriz original.

Esta expansión puede ser hecha para cualquier fila o columna de la matriz original. Por ejemplo, para una expansion en la fila $i$ de una matriz $A$:

\begin{equation}
    D_n = \det(A) = \sum_{j=1}^n a_{ij}(-1)^{i+j}M_{ij}
\end{equation}

Donde $(-1)^{i+j}M_{ij}$ es el cofactor del elemento $a_{ij}$. Luego, para una expansión en la columna $j$, se cambia el indice de la suma por $i$. 

La expansion se simplifica si hay ceros presentes en la columna o fila escogida, ya que no es necesario evaluar sus menores.

\section{Un problema de autovalores}

Sea una matriz $A$,

\begin{equation}
    A = \mqty(2 & 0 & 0 \\ 0 & 1 & 1 \\ 0 & 1 & 1)
\end{equation}

Buscamos resolver el problema

\begin{equation}
    (A - \lambda\mathds{1})\vec{v} = 0
\end{equation}

Como queremos soluciones no triviales, resolvemos la ecuación secular.

\begin{equation}
    \det(A - \lambda\mathds{1}) = \mdet{2-\lambda & 0 & 0 \\ 0 & 1-\lambda & 1 \\ 0 & 1 & 1-\lambda} = 0 
\end{equation}

Expandemos por menores en la primera fila,

\begin{equation}
    \mdet{2-\lambda & 0 & 0 \\ 0 & 1-\lambda & 1 \\ 0 & 1 & 1-\lambda} = -\lambda(2-\lambda)^2 = 0
\end{equation}

y llegamos a que $\lambda_1 = 0, \lambda_2 = 2, \lambda_3 = 2$. Notar que estamos en presencia de un sistema degenerado. Buscamos los autovectores, primero para $\lambda = 0$:

\begin{equation}
    A\vec{v}_1 = 0 \implies 2x = 0, y + z = 0, y + z = 0
\end{equation}

La solución al sistema es $x = 0$ y $z = -y$, por lo tanto el autovector es 

\begin{equation}
    \vec{v}_1 = \mqty(0 \\ 1 \\ -1)
\end{equation}

Ahora, para $\lambda = 2$:

\begin{equation}
    (A - 2\mathds{1})\vec{v} = 0 \implies 0 = 0, -y + z = 0, y - z = 0 
\end{equation}

Y la solución es $x = x$, $y = z$. Como necesitamos romper la degenerancia, podemos hacer uso de que $x$ es un grado de libertad del sistema y escribimos un vector genérico:

\begin{equation}
    \vec{v} = \mqty(x \\ y \\ y)
\end{equation}

Luego, este vector lo podemos reescribir de la siguiente manera:

\begin{equation}
    \vec{v} = x\mqty(1 \\ 0 \\ 0) + y\mqty(0 \\ 1 \\ 1)
\end{equation}

Y de aquí extraemos los dos vectores linealmente independientes que necesitamos.

\begin{equation}
    \vec{v}_2 = \mqty(1 \\ 0 \\ 0) \quad\land\quad \vec{v}_3 = \mqty(0 \\ 1 \\ 1)
\end{equation}

\section{Un problema del Boas - 3.12 P14}

Tenemos el siguiente sistema físico: 

\begin{figure}[h]
    \centering
    \resizebox{0.6\textwidth}{!}{
    \begin{tikzpicture}[
    % Definimos los estilos para las masas y los resortes
    % scale=0.7, transform shape,
    mass/.style={circle, draw, thick, shading=radial, inner color=white!90!gray, outer color=black!60, minimum size=0.8cm},
    spring/.style={thick, decorate, decoration={coil, aspect=0.5, segment length=5pt, amplitude=4pt, pre length=0.6cm, post length=0.6cm}},
]

% ---- Muro Izquierdo ----
\fill [pattern=north east lines] (-0.3,-1) rectangle (0,1);
\draw[thick] (0,-1) -- (0,1);

% ---- Masas ----
% Masa 1
\node[mass] (m1) at (3,0) {};
\node[above=0.4cm] at (m1) {$m$};

% Masa 2
\node[mass] (m2) at (7,0) {};
\node[above=0.4cm] at (m2) {$m$};

% ---- Muro Derecho ----
\fill [pattern=north east lines] (10,-1) rectangle (10.3,1);
\draw[thick] (10,-1) -- (10,1);

% ---- Resortes ----
% Resorte 1 (Muro a Masa 1)
\draw[spring] (0,0) -- node[above=0.4cm] {$k$} (m1.west);

% Resorte 2 (Masa 1 a Masa 2)
\draw[spring] (m1.east) -- node[above=0.4cm] {$2k$} (m2.west);

% Resorte 3 (Masa 2 a Muro)
\draw[spring] (m2.east) -- node[above=0.4cm] {$k$} (10,0);

% ---- Indicadores de desplazamiento (x1 y x2) ----
% Desplazamiento x1
\draw[thick, dashed] (m1.south) -- +(0,-0.8);
\draw[thick, -latex] (3,-0.6) -- (3.8,-0.6) node[right] {$x_1$};

% Desplazamiento x2
\draw[thick, dashed] (m2.south) -- +(0,-0.8);
\draw[thick, -latex] (7,-0.6) -- (7.8,-0.6) node[right] {$x_2$};

\end{tikzpicture}}

\caption{Sistema de dos masas iguales acopladas a tres resortes.}
\end{figure}

Queremos obtener las frecuencias y los modos normales del sistema. La ley de Hooke nos dice que el potencial está dado por:

\begin{equation}
    V = \frac{1}{2}kx^2
\end{equation}

Luego, dada la configuración del sistema, escribimos la expresión para el potencial como:

\begin{equation}
    V = \frac{1}{2}\pqty{kx_1^2 + 2k(x_2-x_1)^2 + kx_2^2}
\end{equation}

Si escribimos las ecuaciones de movimiento para cada masa sabiendo que, por la segunda ley de Newton,

\begin{equation}
    m\ddot{x}_i = -\pdv{V}{x_i} \qc i = 1,2,...
\end{equation}

Tenemos que, para $x_1$,

\begin{equation}
    \begin{aligned}
    m\ddot{x}_1 &= -\pdv{V}{x_1} = -kx_1 + 2k(x_2 - x_1)\\
    \ddot{x}_1 &= \omega_0^2(-3x_1 + 2x_2) \qc \omega_0^2 = \frac{k}{m}
    \end{aligned}
\end{equation}

y para $x_2$:

\begin{equation}
    \begin{aligned}
    m\ddot{x}_2 &= -kx_2 - 2k(x_2 - x_1)\\
    \ddot{x}_2 &= \omega_0^2(2x_1 - 3x_2)
    \end{aligned}
\end{equation}

Proponemos el siguiente ansatz:

\begin{equation}
    x_i(t) = x_0\sin(\omega t)\qc i=1,2,...
\end{equation}

La ecuación diferencial que nos da esta solución es:

\begin{equation}
    \ddot{x}_i = -\omega^2x_i
\end{equation}

Reemplazando lo anterior en nuestras ecuaciones de movimiento nos da el siguiente sistema de ecuaciones:

\begin{equation}
    \begin{cases}
        -\omega^2x_1 = \omega_0^2(-3x_1 + 2x_2) \\
        -\omega^2x_2 = \omega_0^2(2x_1 - 3x_2)
    \end{cases}
\end{equation}

En forma matricial, este sistema se escribe como:

\begin{equation}
    \lambda\mqty(x_1 \\ x_2) = \mqty(3 & -2 \\ -2 & 3) \mqty(x_1 \\ x_2) \qc \lambda=\frac{\omega^2}{\omega_0^2}
\end{equation}

Tenemos un problema de autovalores, por lo que realizamos el procedimiento usual:

\begin{equation}
    \mdet{3-\lambda & -2 \\ -2 & 3-\lambda} = (3-\lambda)^2 - 4 = 0
\end{equation}

Resolviendo para $\lambda$ llegamos a que $\lambda_1 = 5$ y $\lambda_2 = 1$. Luego, por la definición de $\lambda$, vemos que las frecuencias características son

\begin{equation}
    \omega_1 = \sqrt{\frac{5k}{m}} \quad\land\quad \omega_2 = \sqrt{\frac{k}{m}}
\end{equation}

Por último, los modos de vibración los encontramos reemplazando los autovalores. Para $\lambda = 5$, tenemos 

\begin{equation}
    -2x_1 - 2x_2 = 0, -2x_1 + -2x_2 = 0
\end{equation}

con solución $x_2 = -x_1$. Esto nos dice que, cuando $\omega = \omega_1$, las masas se mueven en direcciones opuestas. Luego, para $\lambda = 1$,

\begin{equation}
    2x_1 - 2x_2 = 0, -2x_1 + 2x_2 = 0
\end{equation}

y la solución es $x_1 = x_2$. Tenemos movimiento en conjunto de las masas cuando $\omega = \omega_2$
\end{document}